\begin{abstract}
LLM이 수학 문제를 더 잘 풀게 하는 한 가지 방법은 같은 문제를 여러 번 풀게 하고 그중 맞는 답을 골라내는 것이며, 이를 위해서는 풀이가 서로 충분히 달라야 한다. 최근 \textit{soft prompt} 연구는 학습으로 만든 특수 벡터를 입력에 끼워 넣어 이 다양성을 늘려 왔지만, 효과의 출처가 \emph{학습된 의미} 때문인지 \emph{주입한다는 행위 자체} 때문인지는 가려지지 않았다. 본 논문은 학습 단계를 완전히 빼고 \emph{무작위로 뽑은 벡터}를 입력 끝에 덧붙이는 Random Soft Prompt(RSP)를 분석한다. RSP는 사전학습 임베딩의 통계(평균·분산)에 맞춘 등방 가우시안에서 매 시도마다 새로 뽑힌 벡터로, 학습된 의미가 전혀 없는데도 다섯 개 수학 벤치마크에서 학습된 soft prompt와 비등한 정확도를 낸다. 메커니즘은 두 단계로 작동한다. 처음 보는 무작위 위치를 어텐션이 받아내야 하므로 답의 첫 몇 토큰 분포가 평탄해져 추론 궤적이 갈라지고(\emph{branching}), 답이 길어질수록 그 영향은 자연히 옅어져 응답이 목표로 수렴한다(\emph{goal-directed convergence}). 이 효과는 어텐션의 한 숫자로 정확히 추적되며, 무작위 재추첨이 temperature 같은 기존 샘플링 기법으로는 닿을 수 없는 후보 토큰 영역에까지 양의 확률을 부여함을 이론적으로 보인다. 실험적으로 RSP는 초반 토큰 다양성을 올리고 temperature와 결합하면 Pass@$N$(여러 시도 중 한 번이라도 맞출 확률)을 더 크게 넓히며, 이 효과를 추론에 한정하지 않고 GRPO \citep{shao2024deepseekmath} 학습 단계로 확장해 실용적 적용 가능성을 보인다. 본 연구는 (i) 학습 없는 무작위 주입만으로도 latent reasoning이 작동함을 새롭게 발견하고, (ii) 이론과 실험으로 그 메커니즘을 검증하며, (iii) 추론을 넘어 학습 과정까지 확장한다.
\end{abstract}


